% Reviewable replacement/insertion blocks for ProveIt.
% Pinned commit: 3447bc59b78d138a7f0d536b66ac6e5cc5d5e49f.
% Chapter paths are relative to Analysis/Polylogarithms/docs/manuscript/chapters/.
% Companion report paths are stated separately.
% Do NOT input this entire file: copy each selected block into its stated context.
% Integrate the new proof sections first and preserve/update their labels.

% E1 -- 05-subcritical-stieltjes.tex, subpick:thm:Gaussian-order.
% Replace only the final paragraph beginning "This is a magnitude bound;".
This theorem gives the magnitude comparison. The additional scalar
comparison in Theorem~\ref{newEuler:thm:domination} proves that the same
axis quantity $C(b)=\beta(b)+2^{-b}\eta(b)$ bounds every scaled Euler
remainder, uniformly over $a>0$. That conclusion uses a separate
remainder argument and does not follow from magnitude monotonicity alone.

% E2 -- 05-subcritical-stieltjes.tex, after subpick:thm:axis-maximum.
% Replace the paragraph beginning "The numerical axis experiment finds".
The original numerical axis experiment locates the maximum near
$b_*=1.30221658710124120959$, with height near
$C(b_*)=1.13656110333950956095$. Proposition~\ref{newEuler:prop:narrow}
now certifies rational enclosures for both quantities, including
\[
1.3022165871012412095923717014
<b_*<1.3022165871012412095923717017.
\]
The original uniqueness proof remains independent of the numerical
stationary-point solver. Corollary~\ref{newEuler:cor:sharp} identifies
this maximum as the sharp universal real-order Euler constant.

% E3 -- 05-subcritical-stieltjes.tex, caption of subpick:fig:Gaussian-axis.
% Replace only the caption; keep the existing image and figure label.
\caption{Numerical values of the exact Gaussian axis expression and its
stationary-point marker. The shaded interval is the original proved
bracket $1<b_*<2$; the horizontal line is the sharp critical Euler
constant. Proposition~\ref{newEuler:prop:narrow} supplies independent
exact narrow enclosures for the maximizing point and height.
The additional remainder comparison in
Theorem~\ref{newEuler:thm:domination} proves that this larger Gaussian
maximum is also the sharp universal Euler constant for all positive
real orders.}

% E4 -- 10-discovery.tex, Real-order geometry after the positive-difference theorem.
% Replace the paragraph beginning "The Gaussian experiments have also led to proofs"
% and ending "does not furnish that universal Euler bound.".
The Gaussian experiments have also led to proofs rather than just
narrower searches. Fixed-total monotonicity is proved in
Theorem~\ref{subpick:thm:fixed-total}, and
Theorem~\ref{subpick:thm:axis-maximum} gives the unique global axis
maximum. Its point and height now have exact narrow enclosures in
Proposition~\ref{newEuler:prop:narrow}. Writing
$C(b)=\beta(b)+2^{-b}\eta(b)$ and $C_*=\max_{b>0}C(b)$,
Theorem~\ref{newEuler:thm:domination} proves
\[
0<2^N(E_N-g_{a,b})<C(b)\le C_*
\qquad(a,b>0,\ N\ge1).
\]
Corollary~\ref{newEuler:cor:sharp} proves that $C_*$ is the least
universal constant; the bound is approached as $a\downarrow0$ at
$b=b_*$ and $N=1$. The convenient rational budget $57/50$ is valid
throughout the positive quadrant. The sharper critical/subcritical
constant $\pi/4+\log(2)/2$ in
Corollary~\ref{subpick:cor:sharp-critical} remains available on its
original domain. The number of quartic transition points and the
Bessel continuation remain open; the effective Lerch zero results
are described below.

% M1 -- 07-reflected-moments.tex, immediately after reflected:cor:sharp proof.
% Replace the two-sentence paragraph beginning "This fixed-$m$ corollary does not supply".
This fixed-$m$ corollary retains an error constant that may depend on $m$;
it does not provide its stated relative accuracy when $m$ grows with $n$.
Theorem~\ref{thm:all-ratio} supplies a different approximation, with
relative algebraic error through every fixed order uniformly for
$0\le m\le n$ as $n\to\infty$. Reflection covers the other half of
the nonnegative quadrant. Uniformity of the present exponentially
accurate or least-term approximation for growing $m$ remains separate.

% M2 -- 07-reflected-moments.tex, final status paragraph under
% \paragraph{Verification and remaining questions.}
% Replace from "The fixed-reflected-exponent signed remainder is now supplied"
% through the end of the file; preserve the preceding verification paragraph.
The fixed-reflected-exponent signed remainder is supplied by
Corollary~\ref{reflected:cor:sharp}. Theorem~\ref{thm:all-ratio} now
provides a uniform algebraic expansion when both exponents increase,
including the full range of their ratios. Its phase uses the inverse
coordinate $y=\log\Gamma(x)$ and includes the exponential Jacobian
term. This does not make the endpoint replacement
$B(x)^m\approx\gamma^m x^m$ uniformly valid, nor does it extend the
fixed-$m$ least-term remainder law to growing $m$. Matching that sharper
error scale uniformly remains a further question.

% M3 -- COMPANION REPORT, not a current manuscript chapter input:
% Analysis/Polylogarithms/docs/reports/uniform-transition-continuation/
% sections/04b-proportional-moments.tex.
% Replace the final paragraph headed "What remains open after this theorem.".
\paragraph{What remains open after this theorem.}
The proportional regime on compact positive exponent ratios is resolved,
including uniqueness, curvature and all fixed asymptotic orders.
Theorem~\ref{thm:all-ratio} now supplies a uniform algebraic expansion
joining that regime to the logarithmic transition and continuing over
all nonnegative exponent ratios. It uses a different phase, so its
finite-exponent saddle should not be identified term by term with the
saddle in the present theorem. A proof of the global slope inequality
using only an analytic inequality, without the finite interval
certificate, remains a worthwhile simplification problem. Uniform
exponentially accurate matching and growing-exponent least-term laws
remain open. The original compact-ratio theorem itself retains its
stated hypotheses and error constants.

% M4 -- COMPANION REPORT:
% Analysis/Polylogarithms/docs/reports/uniform-transition-continuation/
% sections/04-joint-moments.tex.
% Replace only the paragraph headed "Editorial correction to the research status.".
\paragraph{Editorial update to the research status.}
The fixed-$m$ statements and their warnings against substituting a
growing reflected exponent remain valid. Theorem~\ref{thm:all-ratio}
now supplies a different uniform expansion through every fixed algebraic
order over all ratios of the two nonnegative exponents. This closes the
uniform algebraic bridge question beyond the compact transition and
proportional regimes. It does not claim that the transition-specific
polynomials in the present theorem have uniform relative remainders for
every unbounded transition parameter. Explicit matching of those
polynomials to the new expansion, and uniform sharp truncation on the
inverse-gamma least-term scale, remain separate questions. These are
extensions of the existing theorems, not corrections to their proofs.

% M5 -- 10-discovery.tex.
% Insert immediately before \subsection{A focused path toward formal verification}.
\subsection{A uniform bridge for reflected log-gamma moments}
Theorem~\ref{thm:all-ratio} gives a complete expansion to every fixed
algebraic order, uniformly for real $0\le m\le n$ as $n\to\infty$.
Reflection covers every pair of nonnegative exponents as their maximum
tends to infinity, including either axis. The proof uses the previously
certified global log-gamma slope inequality and a phase that includes
the exponential Jacobian. It establishes a uniform Gaussian limit on
the same parameter range. Effective remainder constants, uniform
matching to the older transition polynomials, and exponentially accurate
growing-exponent truncation remain useful next problems. The fixed-$m$
residue and least-term theorems retain their distinct stronger error
scales on their stated domains.

% M6 -- COMPANION REPORT 04-joint-moments.tex, Further research questions.
% Replace the first list item beginning "Prove a uniform continuation of the transition when".
\item Determine the maximal noncompact ranges of
$\lambda=me^{-t}$ on which the present transition-specific coefficient
polynomials admit useful uniform relative remainder bounds. Theorem~\ref{thm:all-ratio}
already supplies a different all-ratio algebraic approximation; the
remaining question concerns the validity and matching of this polynomial
representation when $\lambda$ tends to zero or infinity.

% M7 -- Same companion list.
% Replace the item beginning "Build a uniform approximation joining the interior saddle to".
\item Derive explicit uniform matching laws between
Theorem~\ref{thm:all-ratio}, the proportional saddle expansion, and the
transition at $x\asymp1/m$. Quantify their overlap domains and compare
truncation errors without identifying their finite-exponent saddles or
claiming uniform exponential accuracy from an algebraic expansion.

% L1 -- 09-zero-geometry.tex, remark after zeros:thm:expansion.
% Replace the complete remark titled "What uniformity does not say".
\begin{remark}[Scope of the concentration expansion]
The constants in Theorem~\ref{zeros:thm:expansion} may depend on $n$
and on the compact $c$-interval. Its full expansion therefore retains
its fixed-index, compact-scale scope. Theorem~\ref{efflerch:thm:explicit}
provides a different leading zero-location estimate with an explicit
sufficient threshold, uniformly in every root branch and in
$0\le\rho\le1$, without a compact-scale assumption.
Corollary~\ref{efflerch:cor:growing} extends that leading estimate to
a stated growing-index range. It does not extend every coefficient or
remainder in the present $k^{-1}$ expansion uniformly to that range.
\end{remark}

% L2 -- 09-zero-geometry.tex, zeros:thm:zeroasympt.
% Insert after the COMPLETE proof of that theorem, without changing its statement.
For $n\ge2$, Theorem~\ref{efflerch:thm:explicit} gives the sufficient
choice $K_n=576n^2(3n)^{2n-4}$ for saturation, uniformly on
$0\le\rho\le1$, and bounds the logarithmic relative error of every
zero by $8n/\sqrt{k}$. This explicit bound is deliberately conservative.
The fixed-index offsets and higher corrections proved above remain
sharper asymptotic information at fixed $n$; the new estimate controls
the index dependence and the extreme branches.

% L3 -- 10-discovery.tex.
% Replace the complete paragraph under
% \subsection{Zero geometry beyond a fixed Laurent index}.
The complete classical counts are known through Laurent index seven.
The fifth-index transition has threshold three, and
Theorem~\ref{higherzeros:thm:thresholds} proves sharp thresholds four
and five at indices six and seven. These exact small-order results
remain stronger than the general sufficient bound.
Theorem~\ref{efflerch:thm:explicit} now proves saturation at all indices
$n\ge2$ whenever $k\ge576n^2(3n)^{2n-4}$, uniformly in
$0\le\rho\le1$, with simultaneous logarithmic relative error
$8n/\sqrt{k}$ for all zero branches.
Corollary~\ref{efflerch:cor:growing} consequently proves moving-zero
asymptotics for $n\le\alpha\log k/\log\log k$, for each fixed
$0<\alpha<1/2$. The full $n=O(\log k)$ moving-zero expansion remains
open; the existing spectral coefficient theorem does not supply it by
itself. Sharp general thresholds and small-order classification at
higher indices also remain open. Proposition~\ref{efflerch:prop:terminal}
proves strict leftward motion of a simple persistent terminal zero at
arbitrary index, including strict decrease on $[0,1]$ in the saturated
regime. Global monotonicity of all nonterminal branches remains a
separate question; no finite derivative at $\rho=1$ is asserted.

% L4 -- 09-higher-zero-transitions.tex.
% Replace the final paragraph beginning "This finite sequence suggests further threshold questions;".
This finite sequence suggests further threshold questions; it does not
prove a formula in the index. In particular, the assertion
$K_n^{\rm cl}=n-2$ for every $n\ge5$ remains unproved.
Theorem~\ref{efflerch:thm:explicit} supplies an explicit sufficient
saturation threshold for all indices and all $0\le\rho\le1$.
Combining that upper bound with the terminating critical-point method
to determine sharp thresholds remains a concrete next problem.

% C1 -- 10-discovery.tex, append after the last paragraph of research:sec:S8-rejected.
The new weight-eleven and weight-thirteen candidates in
Section~\ref{sec:candidates} extend this collection of frozen research
targets. The retained $S_{10}$ and $S_{12}$ relations are conjectures:
their exact proximity intervals contain zero and prove no identity.
The rejected first $S_{12}$ search vector is a separate record whose
residual interval excludes zero. Neither that rejection nor the earlier
$S_8$ rejection excludes every relation in the corresponding basket.
The $S_6$ mixed reduction remains unresolved. Integral diagonalization
of the single-color shuffle minor does not supply the missing
mixed-color identity or a period-independence theorem.

% A1 -- 10-discovery.tex, append at end of subsection
% \subsection{Certificates, integral structure and arithmetic evaluation}.
Theorem~\ref{oi:thm:integral} additionally gives an integral
diagonalization of the specified odd-weight shuffle minor, determining
its cokernel, Smith factors, ranks over fields, and optimal common
inverse denominator. Its relation to the published double-zeta matrix
and determinant should retain the accompanying attribution.
Corollary~\ref{oi:cor:basket} concerns the quotient by the stated shuffle
rows. Evaluation at a special point, or adjoining further functional
equations, may create additional relations; no numerical period
independence follows from this formal dimension calculation.
